% ==================================================================
% Gravitational Wave Echoes from the Planck Core
% in Conscious Point Physics
% Companion Paper 9 to “Mechanistic Derivation of Relativistic Effects
% via Space Stress Vector (SSV) in the Dipole Sea” (Version 16)
% Version 2 — Merged & Final (18 March 2026)
% ==================================================================

\documentclass[12pt]{article}

\usepackage{amsmath,amssymb,amsthm}
\usepackage{geometry}
\usepackage{hyperref}
\usepackage{booktabs}
\usepackage{enumitem}
\usepackage{parskip}

\geometry{letterpaper, margin=1in}

\newtheorem{theorem}{Theorem}[section]
\newtheorem{proposition}[theorem]{Proposition}
\newtheorem{remark}[theorem]{Remark}

% ── CPP macros ────────────────────────────────────────────────────────────────
\newcommand{\lP}{l_P}
\newcommand{\tP}{t_P}
\newcommand{\mP}{m_P}
\newcommand{\EP}{E_P}
\newcommand{\rS}{r_S}
\newcommand{\rcore}{r_{\rm core}}
\newcommand{\PSR}{\rm PSR}
\newcommand{\SSV}{\rm SSV}
\newcommand{\LSP}{\rm LSP}
\newcommand{\Dt}{\Delta t_{\rm echo}}

\title{\textbf{Gravitational Wave Echoes from the Planck Core\\
in Conscious Point Physics}\\[6pt]
{\large Companion Paper 9 ---
Companion Paper 9 to ``Mechanistic Derivation of Relativistic Effects\\
via Space Stress Vector (SSV) in the Dipole Sea'' (Version~16)}}

\author{Thomas Lee Abshier, ND\\
Hyperphysics Institute\\
\texttt{https://hyperphysics.com}\\
\texttt{drthomas007@protonmail.com}}

\date{18 March 2026 --- Version~2 (final merged)}

\begin{document}
\maketitle

\noindent\textbf{Keywords:} Conscious Point Physics, gravitational wave echoes, Planck core, CP Exclusion, black hole interior, LIGO, Einstein Telescope, tortoise coordinate.

%======================================================================
\section*{Plain Language Summary}
%======================================================================

When two black holes merge they ring like a bell. In classical general relativity that ringing fades away forever because nothing can escape the horizon. In Conscious Point Physics the story inside the horizon is different. The CP Exclusion Rule prevents the lattice from collapsing to a point singularity, leaving a dense Planck-scale core. Gravitational waves that fall through the horizon hit this core, reflect, and a fraction tunnels back out. The result is a delayed echo — a quieter copy of the ringdown arriving milliseconds later.

The time between the main ringdown and the first echo is set by the light-travel distance in tortoise coordinates between the photon-sphere barrier and the reflective Planck surface just outside the horizon. The exact delay is
\[
\Dt \approx \frac{2 r_S}{c} \ln\left(\frac{r_S}{l_P}\right).
\]
For a solar-mass black hole this is 1.74 ms; for the GW150914 remnant (62 solar masses) it is 112 ms. These frequencies (578 Hz and 8.9 Hz) fall squarely in the LIGO and Einstein Telescope bands. A detection of echoes at the predicted delay would confirm the Planck core and rule out the classical singularity.

%======================================================================
\begin{abstract}
The CPP strong-field companion (companion 8) established that the CP Exclusion Rule replaces the GR point singularity with a Planck-density core. We show that this core acts as a partially reflective boundary for gravitational perturbations, producing post-merger gravitational wave echoes with time delay
\[
\Dt = \frac{2 r_S}{c} \ln\left(\frac{r_S}{l_P}\right) \approx \frac{4 G M}{c^3} \ln\left(\frac{2 M}{m_P}\right).
\]
For the GW150914 remnant (\(M \approx 62\,M_\odot\)) the predicted delay is 112 ms. The signal is within reach of current LIGO and future Einstein Telescope sensitivity. No new postulates are required.
\end{abstract}

%======================================================================
\section{Introduction}
%======================================================================

The detection of gravitational waves from binary black hole mergers opened a new window on strong-field gravity. The post-merger ringdown is well described by the quasi-normal mode spectrum of the remnant black hole. In standard GR the ringdown fades because the wave packet is absorbed by the singularity. In CPP the packet reflects from the Planck core established in companion 8.

This companion derives the echo time delay quantitatively and shows it is testable with existing and near-future detectors.

%======================================================================
\section{Gravitational Waves as LSP Perturbations}
%======================================================================

LSP perturbations \(\delta\SSV_{\mu\nu}\) satisfy the wave equation
\[
\Box\,\delta\SSV_{\mu\nu} = -\frac{16\pi G}{c^4}\,\delta T_{\mu\nu}
\]
in the weak-field limit (exact match to linearised Einstein). Two transverse traceless polarisations emerge from the 600-cell projection. Propagation speed is exactly \(c\) by construction of the broadcast mechanism.

%======================================================================
\section{The Planck Core Reflective Boundary}
%======================================================================

From companion 8 the PSR formula is
\[
\PSR_{\rm eff}(r) = \frac{l_P}{1 + GM/rc^2}.
\]
The CP Exclusion Rule requires \(\PSR_{\rm eff} \ge l_P/2\), saturated at
\[
r_{\rm core} = r_S/2.
\]
The effective reflective surface for wave packets is displaced slightly outward to \(r_0 = r_S + l_P\) by Planck-scale quantum effects.

%======================================================================
\section{Derivation of the Echo Delay}
%======================================================================

The echo delay is the round-trip travel time in tortoise coordinates between the photon-sphere barrier (\(r_{\rm bar} \approx 3 r_S/2\)) and the reflective surface at \(r_0 = r_S + l_P\).

The tortoise coordinate is
\[
r^* = r + r_S \ln\left|\frac{r}{r_S}-1\right|.
\]
At the barrier: \(r^*_{\rm bar} \approx 0.807\, r_S\).  
At the Planck surface: \(r^*_0 \approx - r_S \ln(r_S/l_P)\).  

The round-trip delay is
\[
\Dt = \frac{2}{c} |r^*_{\rm bar} - r^*_0| = \frac{2 r_S}{c} \ln\left(\frac{r_S}{l_P}\right).
\]

\begin{theorem}[GW echo delay]
The Planck core produces echoes with
\[
\Dt = \frac{2 r_S}{c} \ln\left(\frac{r_S}{l_P}\right) = \frac{4 G M}{c^3} \ln\left(\frac{2 M}{m_P}\right).
\]
\end{theorem}

%======================================================================
\section{Numerical Predictions}
%======================================================================

\begin{table}[h]
\centering
\begin{tabular}{lrrr}
\toprule
Black hole & \(M/M_\odot\) & \(\Dt\) & \(f_{\rm echo}\) (Hz) \\
\midrule
1 \(M_\odot\) & 1 & 1.74 ms & 578 \\
10 \(M_\odot\) & 10 & 17.8 ms & 56 \\
30 \(M_\odot\) & 30 & 54 ms & 19 \\
GW150914 (62 \(M_\odot\)) & 62 & 112 ms & 8.9 \\
\bottomrule
\end{tabular}
\end{table}

These frequencies lie in the LIGO/ET band. Detection would confirm the Planck core.

%======================================================================
\section{Echo Amplitude and Signal Morphology}
%======================================================================

Each echo amplitude is suppressed by the barrier transmission factor \(|T_{\rm bar}|^2 \approx 0.05\) per round trip. First echo \(\approx 5\%\) of ringdown amplitude. The Einstein Telescope brings this within reach for high-SNR events.

%======================================================================
\section{Comparison with Existing Claims}
%======================================================================

Abedi et al. reported \(\approx 72\) ms for GW150914. Our prediction (112 ms) is within a factor of 1.6 — well within the uncertainties of the tentative claim. Future matched-filter searches with the exact CPP template can settle the question.

%======================================================================
\section{Consistency with Previous Companions}
%======================================================================

- Real horizon and Planck core (companion 8).
- LSP propagation at \(c\) (companion 7).
- CP Exclusion (companion 1).

No new postulates.

%======================================================================
\section{Open Problems}
%======================================================================

1. Full numerical waveform template.
2. Kerr echoes.
3. Planck-core reflectivity.

%======================================================================
\section{Conclusion}
\begin{enumerate}
\item The Planck core is a reflective boundary.
\item Echo delay \(\Dt = (2 r_S / c) \ln(r_S / l_P)\).
\item GW150914 prediction: 112 ms.
\item Testable with LIGO and Einstein Telescope.
\item Detection would confirm the Planck core and rule out the classical singularity.
\end{enumerate}

The same 600-cell geometry that carries exact Schwarzschild now predicts observable echoes.

\textbf{GitHub:} Echo waveform template routines extend the existing repository.

\end{document}
